\documentclass[12pt]{article}

\usepackage{ucs}
\usepackage[utf8x]{inputenc}
\usepackage[russian]{babel}
\usepackage{array,graphicx,dcolumn,multirow,abstract,hanging,fancyhdr,float}

\usepackage{amssymb}
\usepackage{amsmath}
\usepackage{indentfirst}

\newtheorem{theorem}{Теорема}

\DeclareMathOperator{\cond}{cond}
\newcommand{\norm}[1]{\left\| #1 \right\|}

\begin{document}
    \begin{flushright}
		%%.{version}.%%
	\end{flushright}
	\begin{flushright}
		%%.{date}.%%
	\end{flushright}
    \section*{Вычислительный практикум}
    \section*{Практическое занятие \#8.1. Метод простой\\итерации}

	\subsection*{Приведение системы к виду, удобному для итераций}
	
	Для решения системы
	\begin{equation}
		\bold{A}\bold{x}=\bold{b}\label{f1}
	\end{equation}
	с квадратной невырожденной матрицей, методом простой итерации, необходимо предварительно преобразовать эту систему к виду
	\begin{equation}
		\bold{x}=\bold{Bx}+\bold{c},\label{f2}
	\end{equation}
	что есть:
	\begin{equation}
		\begin{gathered}
			x_1=b_{11}x_{1}+b_{12}x_2+b_{13}x_3+\ldots+b_{1n}x_{n}+c_{1},\\
			x_2=b_{21}x_{1}+b_{22}x_2+b_{23}x_3+\ldots+b_{2n}x_{n}+c_{2},\\
			x_3=b_{31}x_{1}+b_{32}x_2+b_{33}x_3+\ldots+b_{3n}x_{n}+c_{3},\\
			\ldots\hspace{65mm}\\
			x_n=b_{n1}x_{1}+b_{n2}x_2+b_{n3}x_3+\ldots+b_{nn}x_{n}+c_{n}.\\
		\end{gathered}
	\end{equation}
	
	Самый простой способ приведения системы (\ref{f1}) к виду (\ref{f2}) состоит в следующем. Из первого уравнения системы (\ref{f1}) выражается неизвестную~$x_1$:
	\begin{equation}
		x_1=a_{11}^{-1}\left(b_1-a_{12}x_2-a_{13}x_3-\ldots-a_{1n}x_n\right),
	\end{equation}
	из второго уравнения --- неизвестную $x_2$:
	\begin{equation}
		x_2=a_{22}^{-1}\left(b_2-a_{21}x_1-a_{23}x_3-\ldots-a_{2n}x_n\right),
	\end{equation}
	и так далее. В результате получаем систему
	\begin{equation}
		\begin{gathered}
			x_1=\hspace{15mm} b_{12}x_2+b_{13}x_3+\ldots+b_{1,n-1}x_{n-1}+b_{1n}x_n+c_1,\\
			x_2=b_{21}x_1\hspace{15mm} +b_{23}x_3+\ldots+b_{2,n-1}x_{n-1}+b_{2n}x_n+c_2,\\
			x_3=b_{31}x_1+b_{32}x_2\hspace{15mm}+\ldots+b_{3,n-1}x_{n-1}+b_{3n}x_n+c_3,\\
			\ldots\hspace{91mm}\\
			x_n=b_{n1}x_1+b_{n2}x_2+b_{n3}x_3+\ldots+ b_{n,n-1}x_{n-1}\hspace{15mm}+c_n,\\
		\end{gathered}
	\end{equation}
	в которой на главной диагонали матрицы $\bold{B}$ находятся нулевые элементы. Остальные элементы вычисляются по формулам
	\begin{equation}
		b_{ij}=-a_{ij}/a_{ii},\quad c_i=b_i/a_{ii},\quad i,j=1,2,\ldots,n,j \neq i.
	\end{equation}
	Очевидно, что для возможности выполнения указанного преобразования необходимо, чтобы диагональные элементы матрицы $\bold{A}$ были ненулевыми.
	
	\subsection*{Алгоритм метода простой итерации}
	
	Выберем некоторое начальное приближение
	\begin{equation}
		\bold{x}^{\left(0\right)}=
		\begin{pmatrix}
			x_1^{\left(0\right)}\\
			x_2^{\left(0\right)}\\
			\vdots\\
			x_n^{\left(0\right)}
		\end{pmatrix}.
	\end{equation}
	Подставляя его в правую часть системы (\ref{f2}) и вычисляя полученное выражение, находим первое приближение
	\begin{equation}
		\bold{x}^{\left(1\right)}=\bold{B}\bold{x}^{\left(0\right)}+\bold{c}.
	\end{equation}
	Подставляя приближение $\bold{x}^{\left(1\right)}$ в правую часть системы (\ref{f2}), получаем
	\begin{equation}
		\bold{x}^{\left(2\right)}=\bold{B}\bold{x}^{\left(1\right)}+\bold{c}.
	\end{equation}
	Продолжая этот процесс далее, получаем последовательность
	\begin{equation}
		\bold{x}^{\left(0\right)}, \bold{x}^{\left(1\right)},\ldots,\bold{x}^{\left(m\right)},\ldots
	\end{equation}
	приближений, вычисляемых по фомруле
	\begin{equation}
		\bold{x}^{\left(k+1\right)}=\bold{B}\bold{x}^{\left(k\right)}+\bold{c},\quad k=0,1,2,\ldots\label{f11}
	\end{equation}
	В развернутой форме записи формула (\ref{f11}) выглядит следующим образом:
	\begin{equation}
		\begin{gathered}
			x_1^{\left(k+1\right)}=b_{11}x_1^{\left(k\right)}+b_{12}x_2^{\left(k\right)}+b_{13}x_3^{\left(k\right)}+\ldots+b_{1n}x_n^{\left(k\right)}+c_1,\\
			x_2^{\left(k+1\right)}=b_{21}x_1^{\left(k\right)}+b_{22}x_2^{\left(k\right)}+b_{23}x_3^{\left(k\right)}+\ldots+b_{2n}x_n^{\left(k\right)}+c_2,\\
			\ldots\quad\quad\quad\quad\quad\quad\quad\quad\quad\quad\quad\quad\quad\quad\quad\quad\quad\quad\quad\quad\\
			x_n^{\left(k+1\right)}=b_{n1}x_1^{\left(k\right)}+b_{n2}x_2^{\left(k\right)}+b_{n3}x_3^{\left(k\right)}+\ldots+b_{nn}x_n^{\left(k\right)}+c_n.\\
		\end{gathered}
	\end{equation}
	
	\subsection*{Сходимость метода простой итерации}
	
	\begin{theorem}[Необходимое условие сходимости]
		Пусть выполнено \linebreak условие
		\begin{equation}
			\left\|\bold{B}\right\|<1.\label{f14}
		\end{equation}
		Тогда
		\begin{enumerate}
			\item решение $\overline{\bold{x}}$, системы (\ref{f2}) существует и единственно;
			\item при произвольном начальном приближении $\bold{x}^{\left(0\right)}$ метод простой итерации сходится и справедлива оценка погрешности
			\begin{equation}
				\left\|\bold{x}^{\left(m\right)}-\overline{\bold{x}}\right\|\leqslant  \left\|\bold{B}\right\|^{m}\left\|\bold{x}^{\left(0\right)}-\overline{\bold{x}}\right\|.
			\end{equation}
		\end{enumerate}
	\end{theorem}
	
	\begin{theorem}
		Если выполнено условие (\ref{f14}), то справедлива апостериорная оценка погрешности
		\begin{equation}
			\left\|\bold{x}^{\left(m\right)}-\overline{\bold{x}}\right\|\leqslant  
			\frac{\left\|\bold{B}\right\|}{1-\left\|\bold{B}\right\|}\left\|\bold{x}^{\left(m\right)}-\bold{x}^{\left(m-1\right)}\right\|.\label{f16}
		\end{equation}
	\end{theorem}
	
	Величину, стоящую в правой части неравенства (\ref{f16}), можно легко вычислить после нахождения очередного приближения $x^{\left(m\right)}$. Если требуется найти решение с точностью $\varepsilon$, то в силу (\ref{f16}) следует вести итерации до выполнения неравенства
	\begin{equation}
		\frac{\left\|\bold{B}\right\|}{1-\left\|\bold{B}\right\|}\left\|\bold{x}^{\left(m\right)}-\bold{x}^{\left(m-1\right)}\right\|<\varepsilon.
	\end{equation}
	
	Таким образом, в качестве критерия окончания итерационного процесса может быть использовано неравенство
	\begin{equation}
		\left\|\bold{x}^{\left(m\right)}-\bold{x}^{\left(m-1\right)}\right\|<\varepsilon_1,\label{f18}
	\end{equation}
	где $$\varepsilon_1=\frac{1-\left\|\bold{B}\right\|}{\left\|\bold{B}\right\|}\varepsilon.$$
	
	На практике иногда используют привлекательный своей простотой критерий окончания
	\begin{equation}
		\left\|\bold{x}^{\left(m\right)}-\bold{x}^{\left(m-1\right)}\right\|<\varepsilon\label{f19}.
	\end{equation}
	Тут стоит отметить, что для метода простой итерации его применение обосновано только тогда, когда $\left\|\bold{B}\right\|\leqslant 1/2$ (в этом случае $$\left(1-\left\|\bold{B}\right\|\right)/\left\|\bold{B}\right\|~\geqslant~1$$
	и выполнение неравенства (\ref{f19}) влечет за собой выполнение неравенства~(\ref{f18})). Однако в большинстве реальных случаев величина $\left\|\bold{B}\right\|$ оказывается близкой к единице и поэтому $$\left(1-\left\|\bold{B}\right\|\right)/\left\|\bold{B}\right\|~\ll~1.$$
	В этих случаях $\varepsilon_1\ll\varepsilon$ и использование критерия (\ref{f19}) приводит к существенно преждевременному окончанию итераций. Величина\linebreak $\left\|\bold{x}^{\left(m\right)}-\bold{x}^{\left(m-1\right)}\right\|$ здесь оказывается малой не потому, что приближения $\bold{x}^{\left(m-1\right)}$ и $\bold{x}^{\left(m\right)}$ близки к решению, а потому, что метод сходится медленно.
	
	Пусть $\lambda_1, \lambda_2,\ldots,\lambda_n$ --- собственные значения матрицы $\bold{B}$. Спектральным радиусом матрицы $\bold{B}$ называется величина $\rho\left(\bold{B}\right)=\max_{1\leqslant i\leqslant n}|\lambda_i|$.
	\begin{theorem}[Необходимое и достаточное условие сходимости]
		\,\\Метод простой итерации (\ref{f11}) сходится при любом начальном приближении $\bold{x}^{\left(0\right)}$ тогда и только тогда, когда $\rho\left(\bold{B}\right)<1$, то есть когда все собственные значения матрицы $\bold{B}$ по модулю меньше единицы.
	\end{theorem}
	
	\subsection*{Задания}
	
	\begin{enumerate}
		\item Привести несколько входных данных, показать их различия.
		\item Выбрать различные начальные приближения --- правая часть, вектор из нулей, вектор из единиц, случайный вектор.
		\item Реализовать метод простой итерации, при вычислениях показать промежуточные результаты каждой итерации. Оценить сходимость метода на каждом шаге и указать сколько итераций потребовалось для достижение некоторой точности.
		\item Требуемая точность приближенного решения $10^{-3}$, $10^{-6}$.
		\item Для прекращения итерационного процесса воспользоваться различными критериями окончания, обосновать их выборы, сравнить их.
		\item Оценить число арифметических операций.
	\end{enumerate}
	
    \renewcommand{\bibname}{{Список литературы}}
    \bibliographystyle{gost2008}
    \begin{thebibliography}{33}
        \bibitem{lebedeva}
        Лебедева А. В., Пакулина А. Н. Практикум по методам вычислений. Часть 1. Учебно-методическое пособие. СПб.: Санкт-Петербургский государственный университет, 2021.~156 с.
        
        \bibitem{Mis}
        Мысовских И. П. Лекции по методам вычислений: Учебное пособие. СПб.: Издательство Санкт-Петербургского университета, 1998. 472~с.
        
        \bibitem{Fadeev}
        Фаддеев Д. К., Фаддеева В. Н. Вычислительные методы линейной алгебры. М.:  Государственное издательство физико-математической литературы, 1960. 655 с.
    \end{thebibliography}
\end{document}
